This project developed and tested an inherited computational fluid dynamics model to investigate two-phase inlet conditions in a geothermal steam--water separator. Separate liquid and vapour inlet regions made the prescribed distribution explicit. Numerical collection, staged loading and wall-film modelling extended the model to examine where liquid accumulated and how it could reach removal. These developments addressed continuing inventory and pressure changes identified during baseline mesh assessment.

On the same mesh and at matched 10,000-iteration endpoints, the imposed stratified inlet retained about 37\% more liquid than the mixed inlet and produced clearer wall and base enrichment. Both still sent nearly all incoming bulk liquid through the steam outlet, so the result established sensitivity of liquid distribution rather than improved separation. Low-feed preparation, inlet ramping and Coupled pressure--velocity treatment supplied starting fields for wall studies, with continuity residuals of approximately 0.002--0.005 at full feed. The reported full-feed state retained excess liquid removal equivalent to 20.81\% of liquid feed.

Wall-film calculations made surface transport, storage and drainage separately observable. Selected windows passed local balance checks with the bulk equations active. A matched direct-drain test reduced the lower-wall store by 70.32\%, but total film by only 1.70\% relative to the undrained case. Local drainage therefore did not control the whole film inventory. Explicit brine-outlet studies added component and short-duration checks, without establishing sustained operating drainage.

Mesh and particle tests identified which predictions remained sensitive to numerical choices. At 195 ms, the finest of four completed frozen-bulk mesh states had about 30\% less film and 32\% less instantaneous drainage than the next coarser mesh. Differences in preparation and held bulk fields prevented an isolated mesh-error estimate. Reconstructed efficiency estimates agreed with earlier project estimates within 0.019 percentage points when unfinished tracks were counted as separated, although 16.78--33.29\% of injected particle mass had no resolved final destination. A tenfold increase in tracking allowance more than doubled the small escaped-mass estimate.

The contribution is a set of tested modelling capabilities and evidence linking inlet distribution, liquid storage and removal. A complete, stationary, mesh-independent and physically validated separator solution was not established. The findings define the liquid accounting, particle-fate and geometry checks still needed before using the model for performance optimisation.
